\section{The polynomial vacuum}\label{sec:vacuum}

The stationary vacuum is initially only a rational vector. To extract a physical observable at coincident parameters, we need a polynomial numerator with a controlled degree. We first construct its scalar normalization and then prove the vector degree bound by interpolation and stationarity.


Exchange equations, degree bounds and reductions at special rapidities
have earlier uses in the dense $O(1)$ model
\cite{DiFrancescoZinnJustin2005} and in the open dilute $O(1)$ model
\cite{GarbaliNienhuis2017Ground,GarbaliNienhuis2017Sum}.
The latter comparison concerns the fusion construction of the polynomial
normalization. A related verification of a Pfaffian normalization by
deletion identities and degree bounds appears for the crossing $O(1)$
model in \cite[Section~2.6]{DiFrancesco2005Open}.
We construct the present zero-loop-value scalar and vector
with their own degree and stationarity proofs.

Write $X=(u_1,\ldots,u_N)$ and $a_i=a(u_i)$. Define
\begin{equation}\label{eq:01-pfaffian}
\begin{split}
 D(a,b)&=a^2+b^2+2dab-d^2,\\
 p_N(a_1,\ldots,a_N)
 &=\prod_{i<j}\frac{D(a_i,a_j)}{a_i-a_j}
   \operatorname{Pf}\left[
    \frac{(a_i-a_j)(a_i+a_j+\rho)}{D(a_i,a_j)}
   \right].
\end{split}
\end{equation}
For odd $N$, augment the skew matrix by a last column of ones and its
negative transpose; set $p_0=1$.

\begin{lemma}[Scalar reduction identities]\label{lem:scalar-reductions}
The function $p_N$ is a symmetric polynomial of degree at most $N-1$
in each argument. Its leading coefficient in any argument is
$p_{N-1}$ of the remaining arguments, and
\begin{equation}\label{eq:01-scalar-flat}
 p_N(Y,c)=p_{N-1}(Y)\prod_{z\in Y}(z+C).
\end{equation}
For $f_j=\cos(w+j\lambda)$,
\begin{align}
 p(Y,f_0,f_6)
 &=p(Y)(f_0+f_6+\rho)
       \prod_{z\in Y}(z-f_{-6})(z-f_{12}),
       \label{eq:01-scalar-delete}\\
 p(Y,f_{-2},f_2)
 &=p(Y,f_0)(f_{-2}+f_2+\rho)
       \prod_{z\in Y}(z+f_0).
       \label{eq:01-scalar-merge}
\end{align}
Here the subscript of $p$ is its number of arguments. Both formulas
also hold with all shifts negated.
\end{lemma}

\begin{proof}
After multiplication by all $D$ factors, the Pfaffian is an alternating
polynomial, hence divisible by the Vandermonde product in
\eqref{eq:01-pfaffian}. Counting degrees before and after division
gives degree at most $N-1$ in each variable. Letting a last argument
tend to infinity proves the leading-coefficient statement: its kernel
entries tend to $-1$. In odd size first add the augmenting column and
row to that argument's column and row. At the flat value,
$D(z,c)=(z-c)(z+C)$ and the kernel entry equals one. This proves
\eqref{eq:01-scalar-flat}.

At $D(f_0,f_6)=0$ only Pfaffian terms pairing these two arguments
survive after denominators are cleared. The identity
$D(f_0,z)=(z-f_{-6})(z-f_6)$ then gives
\eqref{eq:01-scalar-delete}.

Prove the merge identity by induction on $r=|Y|$. The case $r=0$
follows from $p_2(a,b)=a+b+\rho$. If $r\ge1$, distinguish a spectator
$z$. Both sides have degree at most $r+1$ in $z$, and their leading
coefficients agree by induction. They vanish at $z=-f_0$: delete
this argument with $f_{-2}$, and the remaining factor involving
$f_2$ vanishes. At $z=c$ they agree by the flat formula and
\[
 (C+f_{-2})(C+f_2)=(c+f_0)(C+f_0).
\]
They also agree at each deletion value of $z$ with another spectator.
To check the factors, write that pair as $g_0,g_6$, where
$g_j=\cos(v+j\lambda)$, and use
\[
 \prod_{\ell=\pm2}(f_\ell-g_{-6})(f_\ell-g_{12})
 =(f_0+g_0)(f_0+g_6)(f_0-g_{-6})(f_0-g_{12}).
\]
This is the product-to-sum identity obtained from shifts by
$\pm2\lambda$. There are $2(r-1)$ such deletion values, in addition
to $-f_0$ and $c$. For generic arguments these are distinct and give
at least $r+1$ roots for a difference of degree at most $r$, after
its leading coefficient has canceled. Polynomial continuation concludes
the proof.
\end{proof}

\begin{lemma}[Vacuum degree bound]\label{lem:vacuum-degree}
The vector $W_N=p_NP_N$ is Laurent polynomial in $e^{iu_1},\ldots,e^{iu_N}$,
with each exponent between $-(2N-2)$ and $2N-2$. In the last variable
$y=u_N$, $b=a(y)$, its last-slot-vacant part is a polynomial in $b$
of degree at most $N-1$ with leading coefficient $W_{N-1}$.
Its last-slot-occupied part is $\sin y$ times a polynomial in $b$ of
degree at most $N-2$. Set $W_0=1$.
\end{lemma}

\begin{proof}
The assertion for $N=1$ is immediate. We construct the asserted
polynomial by induction and then prove its stationarity. Throughout the
construction, $W$ denotes the candidate at size $N$, whereas every
smaller $W_j=p_jP_j$ is already known by induction. We first remove
the candidate's possible poles and bound its degree, then establish
its special-row relations, and finally use those relations to prove
stationarity.

Deleting a vacant last slot identifies that sector with
$\mathcal L_{N-1}^0$. Deleting an occupied last slot leaves its partner
as the exposed defect of a state in $\mathcal L_{N-1}^1$. We use these
identifications for the subscripts $v$ and $o$ below.

\paragraph{Construction and regularity.}

For each earlier row $i<N$, introduce the deletion node and its
associated cosine values
\[
 y_i=u_i+3\lambda,\quad \beta_i=a(y_i),\quad
 \gamma_i=a(u_i-3\lambda),\quad \sigma_i=a(u_i+6\lambda).
\]
Let $\mathcal R_i$ move the newly inserted row $y_i$ from immediately
above row $i$ to the last position, using successive boxes with
arguments $y_i-u_j$, $i<j<N$. By
\eqref{eq:01-vacuum-exchange}, \eqref{eq:01-vacuum-reduction}, and
\eqref{eq:01-scalar-delete}, the value required at $y=y_i$ is
\begin{equation}\label{eq:01-interpolation-data}
 Y_i=(a_i+\beta_i+\rho)
       \prod_{\substack{j<N\\j\ne i}}
           (a_j-\gamma_i)(a_j-\sigma_i)\,
       \mathcal R_i S_3W_{N-2}(X\setminus\{u_i,u_N\}).
\end{equation}
These data agree with the genuine rational vector $p_NP_N$ on the
node hypersurface: one may check the identity near $u_i=0$ with the
other rows mixed, where all stationary vectors are regular. The mixed point supplies an
open neighborhood within the hypersurface, so this establishes a rational
identity in the remaining parameters, not just equality at one point.

For distinct nodes put
\[
 \ell_i(b)=\prod_{\substack{j<N\\j\ne i}}
                 \frac{b-\beta_j}{\beta_i-\beta_j}.
\]
Define the candidate by the two interpolation formulas
\begin{equation}\label{eq:vacuum-interpolants}
 \begin{aligned}
 W_v(b)&=W_{N-1}\prod_{i<N}(b-\beta_i)
                 +\sum_{i<N}(Y_i)_v\ell_i(b),\\
 W_o(y)&=\sin y\sum_{i<N}\frac{(Y_i)_o}{\sin y_i}\ell_i(b).
 \end{aligned}
\end{equation}
The vacant part has degree at most $N-1$ and leading coefficient
$W_{N-1}$; the occupied part divided by $\sin y$ has degree at most
$N-2$. Both have the prescribed node values. Scalar interpolation
shows that the empty coordinate of $W$ is $p_N$.

We check carefully that this construction has no poles in earlier
spectral variables. Each moving box has two sine factors in its
denominator; they divide the factors
$(a_j-\gamma_i)(a_j-\sigma_i)$ in
\eqref{eq:01-interpolation-data}. The possible divisor $\sin y_i$
in the occupied data also cancels. Explicitly, for $x=u_i$ and
$v=u_j$, the two cancellations are
\[
 \frac{(a(v)-a(x-3\lambda))(a(v)-a(x+6\lambda))}{L(x+3\lambda-v)}
 =-4\sin(x+v-6\lambda)\sin(x+v+3\lambda),
\]
\[
 a(x)+a(x+3\lambda)+\rho
 =4c\sin(x+3\lambda)\sin(3\lambda-x).
\]
The remaining possible divisors come from colliding interpolation
nodes. Since
\[
 \beta_i-\beta_j
 =-2\sin(u_i+u_j+3\lambda)\sin(u_i-u_j),
\]
the two components of $\beta_i=\beta_j$ are
\[
 u_i=u_j\pmod\pi,
 \qquad u_i+u_j=-3\lambda\pmod\pi.
\]
On the first component the data, including the divided occupied data,
agree by their identities with the true vector and its occupancy
signs. This can be verified near the regular mixture
$u_i=u_j=0$, $y_i=y_j=3\lambda$. On the second component both data
vanish through their factors involving $\sigma_i,\sigma_j$; at a
generic point no moving-box denominator vanishes. These checks at
generic points of every irreducible divisor suffice, since all
possible denominators are products of the indicated simple Laurent
factors. Thus $W$ is Laurent polynomial in the earlier variables;
multiple intersections of these divisors introduce no further denominators.

For its degree, send an earlier $u_k$ to either Laurent extreme.
The moving boxes remain bounded there, as follows directly from
the weights in \eqref{eq:01-rhombic-weights}.
Data with $i\ne k$ grow at most like $a_k^{N-1}$: they contain two
scalar factors and the inductive vector $W_{N-2}$, while
$\ell_i(b)$ stays bounded. Data with $i=k$
grow at most like $a_k^{1+2(N-2)}$, while their Lagrange cardinal
polynomial contributes $O(a_k^{-(N-2)})$. The leading-coefficient
term $W_{N-1}\prod_{i<N}(b-\beta_i)$ has the same bound. Since poles
have already been removed, growth at the two Laurent extremes bounds
every earlier exponent between $-(2N-2)$ and $2N-2$. Interpolation preserves
the occupancy sign under a shift $u_k\mapsto u_k+\pi$; when the
shift affects a node, its last-slot sign cancels the sign of the
dividing sine.

\paragraph{Relations in the earlier rows.}
Earlier-row braid covariance follows from the same last-variable
interpolation: the exchange box is independent of the last variable,
the leading coefficients agree by induction, and every node datum agrees
by the rational vacuum exchange identity. Uniqueness of the interpolant
then gives covariance.

We have established the degree bounds and prescribed node values of
$W$. To prove stationarity, we need more zeros than the last-variable
interpolation supplies. The next two row specializations provide those
zeros; the final argument combines them with the degree bounds and
occupancy parity. Both specializations are proved by interpolation in
the last variable.
\begin{enumerate}
\item If adjacent earlier rows satisfy $u_{r+1}=u_r+s\lambda$,
$s=2,3$, the interpolant equals the smaller split vector multiplied
by the corresponding scalar quotient from
\eqref{eq:01-scalar-delete} or \eqref{eq:01-scalar-merge}.
The quotient has degree one in $b$ for merge and two for deletion,
with monic last-variable factor. Thus the required degree forms and leading
coefficient agree by induction. At a node indexed by $i\ne r+1$,
compare both expressions with the true vector near a regular
simultaneous specialization: take $u_r=0$, $u_i=0$ when distinct,
and $y_i=3\lambda$, with the other rows mixed. At $i=r+1$ both
expressions vanish. On the split side the vanishing factor is
$b-\sigma_r$ in deletion or $b+a(u_r+\lambda)$ in merge. On the
node-data side it is $a_r-\gamma_i$ or $a_r-\sigma_i$, respectively;
no moving box crosses row $r$. This proves the specialization.
\item At $u_1=0$ or $3\lambda$, use the first-slot formulas following
\eqref{eq:01-vacuum-reduction}, replacing their smaller vacuum vector
by $W_{N-1}$ and multiplying by $\prod_{j=2}^N(a_j+C)$.
The only degree check requiring attention is the south-source column
in the occupied first-slot part. With north vacant, the last box
contributes $1$ or $A(y)$ to a vacant last output, from vacant or
occupied input respectively, and $U(y)$ or $B(y)$ to an occupied
last output. Multiplication by $b+C$ clears its denominator and gives
exactly the required last-variable forms. Only the contribution $1$
has the leading vacant degree, leaving the formula with the last row
omitted. At the nodes the expressions agree at regular mixed
specializations, taking $u_i=0$ for $i>1$. The exception
$u_1=3\lambda$, $i=1$, $y_i=6\lambda$ is handled by the vanishing
factors $a(y_i)+C$ and $a_1+\beta_1+\rho$.
\end{enumerate}
On each specialization the nodes are generically distinct and their
dividing sines nonzero; all other cases follow by continuation.

\paragraph{Stationarity.}
Consider the polynomial stationarity residual
\[
 V=\left(\prod_{j=1}^N L(u_j)\right)(T_N-I)W.
\]
At each last-variable node it vanishes, because the data came from
the genuine stationary vector. To determine its degree, put $n=N-1$,
$q=W_o/\sin y$ and $\Lambda_n=\prod_{j<N}L(u_j)$. Besides the north-source
map $\mathsf b_n^{\rm north}$, use the map
$\mathsf C_n^{\rm north}:\mathcal L_n^1\to\mathcal L_n^0$ obtained
from a column with north occupied and south vacant by tying its north
end to the exterior defect of the input state. The resulting output
is a vacuum pairing; a diagram that closes a loop contributes zero.
These maps and the transfers on the first $n$ rows are independent of
$y$.

The last box contributes $1$ or $A(y)$ to a vacant output and $U(y)$
or $B(y)$ to an occupied output. Using
$L(y)=(b+C)/2$ and
$\sin y\sin(3\lambda-y)=(b-c)/2$ therefore gives the exact blocks
\begin{align}
 V_v&=\frac{\Lambda_n}2\left((b+C)(T_n-I)W_v
                 +d(b-c)\mathsf C_n^{\rm north}q\right),
                 \label{eq:vacuum-residual-vacant}\\
 \frac{V_o}{\sin y}
   &=\frac{\Lambda_n}2\left((b-c)T_n^{(1)}q-(b+C)q
                 +2d\mathsf b_n^{\rm north}W_v\right).
                 \label{eq:vacuum-residual-occupied}
\end{align}
The second block has degree at most $N-1$. In the first, the only
possible degree-$N$ term has coefficient
\[
 [b^N]V_v=\frac{\Lambda_n}2(T_n-I)W_n=0
\]
by the inductive stationarity. Thus $V_v$ also has degree at most
$N-1$. Consequently every
coordinate of $V$ contains $\prod_{i<N}(b-\beta_i)$.
In each earlier variable $u_k$ its Laurent exponents are bounded by
$2N$. The two sets of specializations just proved make $V$ vanish
when an earlier row is $0$ or $3\lambda$, or when
$u_j=u_i+s\lambda$, $i<j<N$, $s=2,3$. To pass from adjacent to
arbitrary pairs, move the two rows next to one another without
interchanging them, using earlier-row covariance and the column
identity; generic invertibility follows from
\eqref{eq:01-unitarity}.

For fixed $k<N$, these zeros force divisibility by
\[
 (b-\beta_k)\sin u_k\sin(u_k-3\lambda)
 \prod_{\substack{j<k\\s\in\{2,3\}}}
       \sin(u_k-u_j-s\lambda)
 \prod_{\substack{k<j<N\\s\in\{2,3\}}}
       \sin(u_j-u_k-s\lambda).
\]
For generic other variables this divisor has Laurent extremes
$\pm2N$: the factor $b-\beta_k$ contributes two in each direction,
the two flat sine factors contribute two, and the $2(N-2)$ pair factors
contribute the remaining $2N-4$. Its parity is even under $u_k\mapsto u_k+\pi$.
The quotient of a Laurent polynomial with the same exponent bound
by this divisor is independent of $e^{iu_k}$. If slot $k$ is
occupied, however, the corresponding coordinate of $V$ has odd
parity, so that coordinate must be zero. Every nonempty vacuum
pairing has an occupied earlier slot, since it has at least two
occupied slots. The empty residual is identically zero. Hence
$V=0$, so $W$ is stationary. Its empty coordinate is $p_N$;
uniqueness at generic regular parameters gives $W=p_NP_N$.
\end{proof}
